% GENERATED FILE. Edit canonical lessons, then run npm run generate:books.
% canonical-source-hash: fnv1a64:7e48ed416fafda63
% canonical-lessons: FR-C49-bras, FR-C49-jambe, FR-C49-pied, FR-C49-doigt, FR-C49-dent

\chapter{Arm, Leg, Foot, Finger, Tooth}
\label{ch:fr-arm-leg-foot-finger-tooth}
\hlchaptermodality{49}

\begin{chapteropening}
\textbf{By the end of this chapter:} \emph{I can name an arm, a leg, a foot, a finger, a tooth.}

Complete the last lesson of chapter 49: I can name an arm, a leg, a foot, a finger, a tooth.
\end{chapteropening}

\section[le bras]{le bras — an arm}

\label{lesson:FR-C49-bras}



\begin{quote}
Before the new one: say the French for enough, then the French for sir.
\end{quote}

\subsection*{You'll want to know: le bras}
\textbf{le bras} — \textquotedblleft{}an arm\textquotedblright{}.

\begin{cousinweb}
\textbf{bras} $\leftarrow$ Old French \emph{bras} $\leftarrow$ Latin \emph{bracchium}, \textquotedblleft{}arm\textquotedblright{} — inherited. (EtymDB 2.1)
\end{cousinweb}

\begin{grammarlens}[title={What you've built}]
One more word for this chapter.
\end{grammarlens}

\subsection*{Your turn}
\begin{itemize}
\raggedright
  \item \emph{Say it:} \emph{le bras}
  \item \emph{Say it:} \emph{le bras}, once more
  \item \emph{Say it:} say \emph{le bras}
  \item \emph{Recall:} say the French for enough, then the French for sir, then say \emph{le bras} again
\end{itemize}

\begin{tcolorbox}[breakable,colback=teal!4,colframe=teal!35!black,title={Before you move on}]
What does le bras mean? (\textquotedblleft{}An arm\textquotedblright{}.) Say it once more.
\end{tcolorbox}
\section[la jambe]{la jambe — a leg}

\label{lesson:FR-C49-jambe}



\begin{quote}
Before the new one: say the French for sir, then the French for an arm.
\end{quote}

\subsection*{You'll want to know: la jambe}
\textbf{la jambe} — \textquotedblleft{}a leg\textquotedblright{}.

\begin{cousinweb}
\textbf{jambe} $\leftarrow$ Late Latin \emph{gamba}, \textquotedblleft{}a hoof, a leg\textquotedblright{} (Lewis \& Short) — inherited: Latin \emph{g} before \emph{a} softened to the French \emph{j}. EtymDB 2.1 traces Italian \emph{gamba} to the same Late Latin word.
\end{cousinweb}

\begin{grammarlens}[title={What you've built}]
One more word for this chapter.
\end{grammarlens}

\subsection*{Your turn}
\begin{itemize}
\raggedright
  \item \emph{Say it:} \emph{la jambe}
  \item \emph{Say it:} \emph{la jambe}, once more
  \item \emph{Say it:} say \emph{la jambe}
  \item \emph{Recall:} say the French for sir, then the French for an arm, then say \emph{la jambe} again
\end{itemize}

\begin{tcolorbox}[breakable,colback=teal!4,colframe=teal!35!black,title={Before you move on}]
What does la jambe mean? (\textquotedblleft{}A leg\textquotedblright{}.) Say it once more.
\end{tcolorbox}
\section[le pied]{le pied — a foot}

\label{lesson:FR-C49-pied}



\begin{quote}
Before the new one: say the French for an arm, then the French for a leg.
\end{quote}

\subsection*{You'll want to know: le pied}
\textbf{le pied} — \textquotedblleft{}a foot\textquotedblright{}.

\begin{cousinweb}
\textbf{pied} $\leftarrow$ Old French \emph{pié} $\leftarrow$ Latin \emph{pes}, \emph{pedis}, \textquotedblleft{}foot\textquotedblright{} — inherited. (EtymDB 2.1)
\end{cousinweb}

\begin{grammarlens}[title={What you've built}]
One more word for this chapter.
\end{grammarlens}

\subsection*{Your turn}
\begin{itemize}
\raggedright
  \item \emph{Say it:} \emph{le pied}
  \item \emph{Say it:} \emph{le pied}, once more
  \item \emph{Say it:} say \emph{le pied}
  \item \emph{Recall:} say the French for an arm, then the French for a leg, then say \emph{le pied} again
\end{itemize}

\begin{tcolorbox}[breakable,colback=teal!4,colframe=teal!35!black,title={Before you move on}]
What does le pied mean? (\textquotedblleft{}A foot\textquotedblright{}.) Say it once more.
\end{tcolorbox}
\section[le doigt]{le doigt — a finger}

\label{lesson:FR-C49-doigt}



\begin{quote}
Before the new one: say the French for a leg, then the French for a foot.
\end{quote}

\subsection*{You'll want to know: le doigt}
\textbf{le doigt} — \textquotedblleft{}a finger\textquotedblright{}.

\begin{cousinweb}
\textbf{doigt} $\leftarrow$ Old French \emph{doit} $\leftarrow$ Latin \emph{digitus}, \textquotedblleft{}finger\textquotedblright{} — inherited. (EtymDB 2.1) The silent \emph{g} was put back into the spelling later, to recall the Latin.
\end{cousinweb}

\begin{grammarlens}[title={What you've built}]
One more word for this chapter.
\end{grammarlens}

\subsection*{Your turn}
\begin{itemize}
\raggedright
  \item \emph{Say it:} \emph{le doigt}
  \item \emph{Say it:} \emph{le doigt}, once more
  \item \emph{Say it:} say \emph{le doigt}
  \item \emph{Recall:} say the French for a leg, then the French for a foot, then say \emph{le doigt} again
\end{itemize}

\begin{tcolorbox}[breakable,colback=teal!4,colframe=teal!35!black,title={Before you move on}]
What does le doigt mean? (\textquotedblleft{}A finger\textquotedblright{}.) Say it once more.
\end{tcolorbox}
\section[la dent]{la dent — a tooth}

\label{lesson:FR-C49-dent}



\begin{quote}
Before the new one: say the French for a foot, then the French for a finger.
\end{quote}

\subsection*{You'll want to know: la dent}
\textbf{la dent} — \textquotedblleft{}a tooth\textquotedblright{}. English \emph{dentist} shares the root.

\begin{cousinweb}
\textbf{dent} $\leftarrow$ Latin \emph{dens}, \emph{dentis}, \textquotedblleft{}tooth\textquotedblright{} — inherited. (EtymDB 2.1)
\end{cousinweb}

\begin{grammarlens}[title={What you've built}]
One more word for this chapter.
\end{grammarlens}

\subsection*{Your turn}
\begin{itemize}
\raggedright
  \item \emph{Say it:} \emph{la dent}
  \item \emph{Say it:} \emph{la dent}, once more
  \item \emph{Say it:} say \emph{la dent}
  \item \emph{Recall:} say the French for a foot, then the French for a finger, then say \emph{la dent} again
\end{itemize}

\begin{tcolorbox}[breakable,colback=teal!4,colframe=teal!35!black,title={Before you move on}]
What does la dent mean? (\textquotedblleft{}A tooth\textquotedblright{}.) Say it once more.
\end{tcolorbox}
